
\documentclass[12pt]{article}
%%%%%%%%%%%%%%%%%%%%%%%%%%%%%%%%%%%%%%%%%%%%%%%%%%%%%%%%%%%%%%%%%%%%%%%%%%%%%%%%%%%%%%%%%%%%%%%%%%%%%%%%%%%%%%%%%%%%%%%%%%%%%%%%%%%%%%%%%%%%%%%%%%%%%%%%%%%%%%%%%%%%%%%%%%%%%%%%%%%%%%%%%%%%%%%%%%%%%%%%%%%%%%%%%%%%%%%%%%%%%%%%%%%%%%%%%%%%%%%%%%%%%%%%%%%%
\usepackage{amsmath}
\usepackage[compat2]{geometry}

\setcounter{MaxMatrixCols}{10}
%TCIDATA{TCIstyle=LaTeX article (bright).cst}

%TCIDATA{OutputFilter=LATEX.DLL}
%TCIDATA{Version=5.50.0.2953}
%TCIDATA{<META NAME="SaveForMode" CONTENT="1">}
%TCIDATA{BibliographyScheme=Manual}
%TCIDATA{Created=Tuesday, August 14, 2012 13:45:33}
%TCIDATA{LastRevised=Monday, September 27, 2021 11:09:32}
%TCIDATA{<META NAME="GraphicsSave" CONTENT="32">}
%TCIDATA{<META NAME="DocumentShell" CONTENT="Exams and Syllabi\SW\Assignment">}

\setlength{\topmargin}{-1.0in}
\setlength{\textheight}{9.25in}
\setlength{\oddsidemargin}{0.0in}
\setlength{\evensidemargin}{0.0in}
\setlength{\textwidth}{6.5in}
\def\labelenumi{\arabic{enumi}.}
\def\theenumi{\arabic{enumi}}
\def\labelenumii{(\alph{enumii})}
\def\theenumii{\alph{enumii}}
\def\p@enumii{\theenumi.}
\def\labelenumiii{\arabic{enumiii}.}
\def\theenumiii{\arabic{enumiii}}
\def\p@enumiii{(\theenumi)(\theenumii)}
\def\labelenumiv{\arabic{enumiv}.}
\def\theenumiv{\arabic{enumiv}}
\def\p@enumiv{\p@enumiii.\theenumiii}
\pagestyle{plain}
\setcounter{secnumdepth}{0}
\parindent=0pt
\input{tcilatex}
\begin{document}


\begin{center}
\begin{equation*}
\FRAME{itbpF}{1.3214in}{0.9496in}{0in}{}{}{Figure}{\special{language
"Scientific Word";type "GRAPHIC";maintain-aspect-ratio TRUE;display
"USEDEF";valid_file "T";width 1.3214in;height 0.9496in;depth
0in;original-width 1.3552in;original-height 0.9669in;cropleft "0";croptop
"1";cropright "1";cropbottom "0";tempfilename
'R03JBW00.bmp';tempfile-properties "XPR";}}
\end{equation*}%
\textbf{Electromagnetismo - 230045}

\emph{Gu\'{\i}a N}$%
%TCIMACRO{\U{b0}}%
%BeginExpansion
{{}^\circ}%
%EndExpansion
4$\emph{\ - Campo el\'{e}ctrico de distribuciones continuas}

\textit{Profesor: Arturo Fern\'{a}ndez P\'{e}rez}
\end{center}

\begin{enumerate}
\item Un conductor en forma de anillo, cuyo radio es $R=7$ $[cm]$, tiene una
carga total $Q=4,5\times 10^{-4}$ [C] distribu\'{\i}da uniformemente en toda
su circunferencia. Determine la fuerza el\'{e}ctrica ejercida sobre una
carga puntual $q=-5$ [$\mu C$] situada sobre el eje del anillo en el punto
P, a una distancia $z=15$ [cm] del centro del anillo, como muestra la Fig.
1. \textbf{R: }$\vec{F}=-669,6$ $\hat{k}$ $[N]$%
\begin{equation*}
\FRAME{itbpFU}{0.9945in}{1.7244in}{0in}{\Qcb{Figura 1}}{}{Figure}{\special%
{language "Scientific Word";type "GRAPHIC";maintain-aspect-ratio
TRUE;display "USEDEF";valid_file "T";width 0.9945in;height 1.7244in;depth
0in;original-width 1.5627in;original-height 2.7294in;cropleft "0";croptop
"1";cropright "1";cropbottom "0";tempfilename
'R03JBW01.wmf';tempfile-properties "XPR";}}
\end{equation*}

\item Una l\'{\i}nea de carga con densidad uniforme $\lambda =4$ $[nC/m]$ se
encuentra sobre el eje $y$, entre $y=0$ e $y=30$ [cm]. Si adem\'{a}s hay una
carga puntual negativa sobre el eje $x$, $q=-5$ [$mC$], a una distancia de
50 cm del origen, como muestra la Fig. 2, determine el campo el\'{e}ctrico
total sobre el eje $x,$ en la posici\'{o}n $x=15$ [cm]. \textbf{R: }$\vec{E}%
=3,673\times 10^{8}\hat{\imath}-132,696$ $\hat{\jmath}$ $[N/C]$ 
\begin{equation*}
\FRAME{itbpFU}{1.7841in}{1.3569in}{0in}{\Qcb{Figura 2}}{}{Figure}{\special%
{language "Scientific Word";type "GRAPHIC";maintain-aspect-ratio
TRUE;display "USEDEF";valid_file "T";width 1.7841in;height 1.3569in;depth
0in;original-width 6.9272in;original-height 5.2641in;cropleft "0";croptop
"1";cropright "1";cropbottom "0";tempfilename
'R03JBW02.wmf';tempfile-properties "XPR";}}
\end{equation*}

\item Una l\'{\i}nea de carga, con distribuci\'{o}n lineal constante $%
\lambda =4\times 10^{-6}$ [$\frac{C}{m}$], se encuentra paralela al eje $x,$
entre \thinspace $x=-1$ $[m]$ y $x=3$ $[m]$, en la posici\'{o}n $z=2$ $[m]$
e $y=0$ $[m]$. Adem\'{a}s, hay una carga puntual $q=-3\times 10^{-6}$ $[C]\,$%
\ en el punto $(4,4,3)$ [m].

\begin{enumerate}
\item Dibuje(mos) un esquema de la situaci\'{o}n.

\item Determine el campo el\'{e}ctrico en el punto $P=(1,6,1)$ [m], inducido
por la carga puntual \ \textbf{R: }$\vec{E}=-1155,611\hat{\imath}+770,407%
\hat{\jmath}-770,407\hat{k}$ [N/C]

\item Determine el campo el\'{e}ctrico en el punto $P=(1,6,1)$ [m], inducido
por la distribuci\'{o}n lineal \textbf{R: }$\vec{E}=3672\hat{\jmath}-612\hat{%
k}$ [N/C]

\item El campo el\'{e}ctrico total en el punto $P=(1,6,1)$ [m] \textbf{R: }$%
\vec{E}=-1155,611\hat{\imath}+4442,407\hat{\jmath}-1382,407\hat{k}$ [N/C]

\item Si en ese punto se coloca una carga puntual~$q_{2}=5$ [mC], determine
la fuerza neta ejercida sobre ella. \textbf{R: }$\vec{F}=-5,778\hat{\imath}%
+22,212\hat{\jmath}-6,912\hat{k}$ [N]
\end{enumerate}

\item Un semic\'{\i}rculo de radio $a=45$ [cm] se encuentra en el primer y
segundo cuadrante, como muestra la Fig. 3. Hay una carga total distribu\'{\i}%
da en el semianillo igual a cero, de manera que hay una carga $-Q=-45$ [mC]
en el primer cuadrante y una carga $+Q=45$ [mC] en el segundo cuadrante.
Determine el campo el\'{e}ctrico total en el origen. \textbf{R: }$\vec{E}%
=2,56\times 10^{9}\hat{\imath}$ [N/C] \textbf{\ } 
\begin{equation*}
\FRAME{itbpFU}{2.2061in}{1.5177in}{0in}{\Qcb{Figura 3}}{}{Figure}{\special%
{language "Scientific Word";type "GRAPHIC";maintain-aspect-ratio
TRUE;display "USEDEF";valid_file "T";width 2.2061in;height 1.5177in;depth
0in;original-width 4.8326in;original-height 3.3148in;cropleft "0";croptop
"1";cropright "1";cropbottom "0";tempfilename
'R03JBW03.wmf';tempfile-properties "XPR";}}
\end{equation*}

\item Considere un disco de radio $R=50$ [cm] con una densidad de carga $%
\sigma =8r^{2}$ [$\frac{\mu C}{m^{2}}$]. Determine el campo el\'{e}ctrico en
el punto $P,$ubicado en el eje del disco a una distancia $z=80$ [cm] del
centro del disco, como muestra la Fig. 4.\textbf{\ R: }$\vec{E}=7,888\times
10^{9}\hat{k}$ [N/C] 
\begin{equation*}
\FRAME{itbpFU}{0.9868in}{1.9545in}{0in}{\Qcb{Figura 4}}{}{Figure}{\special%
{language "Scientific Word";type "GRAPHIC";maintain-aspect-ratio
TRUE;display "USEDEF";valid_file "T";width 0.9868in;height 1.9545in;depth
0in;original-width 1.9899in;original-height 3.9686in;cropleft "0";croptop
"1";cropright "1";cropbottom "0";tempfilename
'R03JBW04.wmf';tempfile-properties "XPR";}}
\end{equation*}%
\begin{equation*}
\text{Indicaci\'{o}n:}\int \frac{r^{3}}{\left( r^{2}+a^{2}\right) ^{3/2}}dr=%
\frac{2a^{2}+r^{2}}{\sqrt{a^{2}+r^{2}}}
\end{equation*}
\end{enumerate}

\end{document}
